%% GENERATED by thesis/tools/gen_supp.py -- do not edit
\begin{table}[!t]
  \centering
  \caption{외부 표본의 짝 경기 군집 부트스트랩 (16.13)}
  \label{tab:ext-bootstrap}
  \footnotesize

  \begin{tabular}{@{}llrrrrrll@{}}
    \toprule
    코호트 & 표본 & $n$ & 경기 & Brier $q$ & Brier PT & $\Delta$Brier & 95\% 구간 & 98.75\% 구간 \\
    \midrule
    T & KR 16.13 & 5{,}202 & 3{,}859 & 0.2483 & 0.2457 & +0.00261 & [-0.00041, +0.00572] & [-0.00115, +0.00668] \\
    T & NA1 16.13 & 5{,}312 & 3{,}955 & 0.2534 & 0.2494 & +0.00398 & [+0.00086, +0.00705] & [+0.00009, +0.00784] \\
    S & KR 16.13 & 15{,}641 & 7{,}874 & 0.2494 & 0.2512 & -0.00180 & [-0.00328, -0.00029] & [-0.00373, +0.00011] \\
    S & NA1 16.13 & 16{,}100 & 7{,}952 & 0.2503 & 0.2523 & -0.00204 & [-0.00363, -0.00044] & [-0.00403, +0.00001] \\
    \bottomrule
  \end{tabular}
  \tabnote{동결 $\hat V$, $q$, $\mathrm{PT}_{\mathrm{flex}}$(S는 $q_S$, $\mathrm{PT}_{\mathrm{flex}}^S$)의 점수; 재학습·재보정 없음. Brier는 경기 가중, $\Delta$Brier는 경기 평균 차이($q - \mathrm{PT}$, 음수이면 $q$ 우세). 구간은 10{,}000회 경기 군집 백분위 부트스트랩(시드 7); 98.75\%는 네 대비 Bonferroni. 구간은 고정된 모델에 조건부이며 모델 적합 불확실성을 포함하지 않는다.}
\end{table}
